% positivstellensatz.tex — Putinar's Positivstellensatz and the constrained
% Lyapunov certificate built on it.
%
% THIS IS THE FLAGSHIP DOCUMENT of the repository. It is the result cited in the
% TVF paper, and it is the one an advisor at a theory-heavy lab will actually
% read. Budget accordingly: it should be the longest and most careful of the two
% derivations, and Section 5 (the ROA theorem as this repository states it) must
% be airtight.
%
% Depends on sos_decomposition.tex for the Gram parameterisation. Cite it; do not
% restate it.
%
% Compile:  latexmk -pdf positivstellensatz.tex   (from docs/derivations/)

\documentclass[11pt]{article}
% preamble.tex — shared macros for the derivations.
%
% Both derivation files are standalone-compilable via % preamble.tex — shared macros for the derivations.
%
% Both derivation files are standalone-compilable via % preamble.tex — shared macros for the derivations.
%
% Both derivation files are standalone-compilable via \input{preamble}. Compile
% from inside docs/derivations/ so the relative \input resolves:
%
%   latexmk -pdf sos_decomposition.tex
%
% Notation here must match docs/theory_notes.md §1 and the code identifiers.
% If they drift, the cross-references in the notebooks stop meaning anything.

\usepackage[utf8]{inputenc}
\usepackage[T1]{fontenc}
\usepackage{amsmath, amssymb, amsthm}
\usepackage{mathtools}
\usepackage{bm}
\usepackage{graphicx}
\usepackage{booktabs}
\usepackage[margin=1in]{geometry}
\usepackage[colorlinks=true, linkcolor=blue, citecolor=blue]{hyperref}
\usepackage{cleveref}   % requires two pdflatex passes

\theoremstyle{plain}
\newtheorem{theorem}{Theorem}
\newtheorem{lemma}[theorem]{Lemma}
\newtheorem{proposition}[theorem]{Proposition}
\newtheorem{corollary}[theorem]{Corollary}

\theoremstyle{definition}
\newtheorem{definition}[theorem]{Definition}
\newtheorem{example}[theorem]{Example}
\newtheorem{assumption}[theorem]{Assumption}

\theoremstyle{remark}
\newtheorem{remark}[theorem]{Remark}

% --- sets and cones ----------------------------------------------------------
\newcommand{\R}{\mathbb{R}}
\newcommand{\N}{\mathbb{N}}
\newcommand{\Sym}{\mathbb{S}}                      % symmetric matrices
\newcommand{\PSD}{\succeq 0}
\newcommand{\SOS}{\Sigma}                          % the SOS cone, \SOS[x]
\newcommand{\Rx}{\R[x]}                            % polynomial ring
\newcommand{\Mset}{\mathcal{M}}                    % quadratic module

% --- the objects of the certificate ------------------------------------------
\newcommand{\Vdot}{\dot{V}}
\newcommand{\Lie}[2]{\nabla #1 \cdot #2}           % \Lie{V}{f}
\newcommand{\sublevel}[2]{\Omega_{#2}}             % \sublevel{V}{\rho}
\newcommand{\ROA}{\mathcal{R}}                     % region of attraction
\newcommand{\gram}{Q}
\newcommand{\basis}{z}

% --- operators ---------------------------------------------------------------
\DeclareMathOperator{\tr}{tr}
\DeclareMathOperator{\rank}{rank}
\DeclareMathOperator{\diag}{diag}
\DeclareMathOperator{\deg}{deg}
\DeclareMathOperator*{\argmax}{arg\,max}
\DeclareMathOperator*{\argmin}{arg\,min}

% --- a macro for pointing prose at the implementation -------------------------
% Use \codelink{lie_derivative}{julia/src/polynomial_utils.jl} in the text every
% time a symbol first appears. A derivation that names the function implementing
% it is checkable; one that does not is a separate document that happens to sit
% in the same repository.
\newcommand{\codelink}[2]{\texttt{#1} (\texttt{#2})}
. Compile
% from inside docs/derivations/ so the relative \input resolves:
%
%   latexmk -pdf sos_decomposition.tex
%
% Notation here must match docs/theory_notes.md §1 and the code identifiers.
% If they drift, the cross-references in the notebooks stop meaning anything.

\usepackage[utf8]{inputenc}
\usepackage[T1]{fontenc}
\usepackage{amsmath, amssymb, amsthm}
\usepackage{mathtools}
\usepackage{bm}
\usepackage{graphicx}
\usepackage{booktabs}
\usepackage[margin=1in]{geometry}
\usepackage[colorlinks=true, linkcolor=blue, citecolor=blue]{hyperref}
\usepackage{cleveref}   % requires two pdflatex passes

\theoremstyle{plain}
\newtheorem{theorem}{Theorem}
\newtheorem{lemma}[theorem]{Lemma}
\newtheorem{proposition}[theorem]{Proposition}
\newtheorem{corollary}[theorem]{Corollary}

\theoremstyle{definition}
\newtheorem{definition}[theorem]{Definition}
\newtheorem{example}[theorem]{Example}
\newtheorem{assumption}[theorem]{Assumption}

\theoremstyle{remark}
\newtheorem{remark}[theorem]{Remark}

% --- sets and cones ----------------------------------------------------------
\newcommand{\R}{\mathbb{R}}
\newcommand{\N}{\mathbb{N}}
\newcommand{\Sym}{\mathbb{S}}                      % symmetric matrices
\newcommand{\PSD}{\succeq 0}
\newcommand{\SOS}{\Sigma}                          % the SOS cone, \SOS[x]
\newcommand{\Rx}{\R[x]}                            % polynomial ring
\newcommand{\Mset}{\mathcal{M}}                    % quadratic module

% --- the objects of the certificate ------------------------------------------
\newcommand{\Vdot}{\dot{V}}
\newcommand{\Lie}[2]{\nabla #1 \cdot #2}           % \Lie{V}{f}
\newcommand{\sublevel}[2]{\Omega_{#2}}             % \sublevel{V}{\rho}
\newcommand{\ROA}{\mathcal{R}}                     % region of attraction
\newcommand{\gram}{Q}
\newcommand{\basis}{z}

% --- operators ---------------------------------------------------------------
\DeclareMathOperator{\tr}{tr}
\DeclareMathOperator{\rank}{rank}
\DeclareMathOperator{\diag}{diag}
\DeclareMathOperator{\deg}{deg}
\DeclareMathOperator*{\argmax}{arg\,max}
\DeclareMathOperator*{\argmin}{arg\,min}

% --- a macro for pointing prose at the implementation -------------------------
% Use \codelink{lie_derivative}{julia/src/polynomial_utils.jl} in the text every
% time a symbol first appears. A derivation that names the function implementing
% it is checkable; one that does not is a separate document that happens to sit
% in the same repository.
\newcommand{\codelink}[2]{\texttt{#1} (\texttt{#2})}
. Compile
% from inside docs/derivations/ so the relative \input resolves:
%
%   latexmk -pdf sos_decomposition.tex
%
% Notation here must match docs/theory_notes.md §1 and the code identifiers.
% If they drift, the cross-references in the notebooks stop meaning anything.

\usepackage[utf8]{inputenc}
\usepackage[T1]{fontenc}
\usepackage{amsmath, amssymb, amsthm}
\usepackage{mathtools}
\usepackage{bm}
\usepackage{graphicx}
\usepackage{booktabs}
\usepackage[margin=1in]{geometry}
\usepackage[colorlinks=true, linkcolor=blue, citecolor=blue]{hyperref}
\usepackage{cleveref}   % requires two pdflatex passes

\theoremstyle{plain}
\newtheorem{theorem}{Theorem}
\newtheorem{lemma}[theorem]{Lemma}
\newtheorem{proposition}[theorem]{Proposition}
\newtheorem{corollary}[theorem]{Corollary}

\theoremstyle{definition}
\newtheorem{definition}[theorem]{Definition}
\newtheorem{example}[theorem]{Example}
\newtheorem{assumption}[theorem]{Assumption}

\theoremstyle{remark}
\newtheorem{remark}[theorem]{Remark}

% --- sets and cones ----------------------------------------------------------
\newcommand{\R}{\mathbb{R}}
\newcommand{\N}{\mathbb{N}}
\newcommand{\Sym}{\mathbb{S}}                      % symmetric matrices
\newcommand{\PSD}{\succeq 0}
\newcommand{\SOS}{\Sigma}                          % the SOS cone, \SOS[x]
\newcommand{\Rx}{\R[x]}                            % polynomial ring
\newcommand{\Mset}{\mathcal{M}}                    % quadratic module

% --- the objects of the certificate ------------------------------------------
\newcommand{\Vdot}{\dot{V}}
\newcommand{\Lie}[2]{\nabla #1 \cdot #2}           % \Lie{V}{f}
\newcommand{\sublevel}[2]{\Omega_{#2}}             % \sublevel{V}{\rho}
\newcommand{\ROA}{\mathcal{R}}                     % region of attraction
\newcommand{\gram}{Q}
\newcommand{\basis}{z}

% --- operators ---------------------------------------------------------------
\DeclareMathOperator{\tr}{tr}
\DeclareMathOperator{\rank}{rank}
\DeclareMathOperator{\diag}{diag}
\DeclareMathOperator{\deg}{deg}
\DeclareMathOperator*{\argmax}{arg\,max}
\DeclareMathOperator*{\argmin}{arg\,min}

% --- a macro for pointing prose at the implementation -------------------------
% Use \codelink{lie_derivative}{julia/src/polynomial_utils.jl} in the text every
% time a symbol first appears. A derivation that names the function implementing
% it is checkable; one that does not is a separate document that happens to sit
% in the same repository.
\newcommand{\codelink}[2]{\texttt{#1} (\texttt{#2})}


\title{Putinar's Positivstellensatz and Certified Regions of Attraction}
\author{Ali-Eimaan}
\date{}

\begin{document}
\maketitle

\begin{abstract}
The S-procedure used throughout this toolbox is an instance of Putinar's
Positivstellensatz: a representation theorem that says every polynomial
strictly positive on a compact basic closed semialgebraic set belongs to the
associated quadratic module. This note states the theorem with its
Archimedean hypothesis, derives the region-of-attraction SOS program as a
special case, and is explicit about which implications are one-way: SOS
feasibility implies invariance, but infeasibility proves nothing about the
system. The final section connects the certified sets produced here to the
transition-viable funnel composition that the broader research programme
requires.
\end{abstract}

\section{Semialgebraic sets and quadratic modules}
Let $\R[x] = \R[x_1,\dots,x_n]$ and let $\Sigma[x] \subset \R[x]$ denote
the cone of sum-of-squares (SOS) polynomials. A \emph{basic closed
semialgebraic set} is
\begin{equation}
  K = \{\, x \in \R^n : g_1(x) \ge 0,\; \dots,\; g_m(x) \ge 0 \,\},
\end{equation}
where $g_j \in \R[x]$. The \emph{quadratic module} generated by
$g = (g_1,\dots,g_m)$ is
\begin{equation}
  M(g) = \Bigl\{\, \sigma_0 + \sum_{j=1}^m \sigma_j g_j
    : \sigma_0,\dots,\sigma_m \in \Sigma[x] \,\Bigr\}.
\end{equation}
The \emph{preordering} $P(g)$ replaces the single sum by a sum over all
$2^m$ products $\prod_{j \in J} g_j$ with $J \subseteq \{1,\dots,m\}$.
Putinar's theorem uses $M(g)$, not $P(g)$, because the quadratic module is
linear in the generators while the preordering introduces $2^m$ terms --- a
practical difference that makes Putinar's formulation the one implemented in
every SOS toolbox \cite{Putinar1993, Parrilo2000}.

\section{The Archimedean condition}
\begin{definition}[Archimedean]
  $M(g)$ is \emph{Archimedean} if there exists $N \in \R$ such that
  $N - \|x\|^2 \in M(g)$.
\end{definition}
The Archimedean condition implies $K$ is compact, but is strictly stronger
than compactness: the distinction is representation-dependent.

The practical consequence is that any compact $K$ can be made to satisfy
the hypothesis by adding the redundant ball constraint $N - \|x\|^2 \ge 0$
to its description --- this leaves $K$ unchanged but makes the quadratic
module Archimedean. The toolbox does exactly this: in
\codelink{add\_ball\_constraint}{julia/src/sos\_lyapunov.jl}, a constraint
$R^2 - \|x\|^2 \ge 0$ is appended with $R$ chosen from the bounding box of
the level set. For the unicycle example, $R = 5$.

\section{Putinar's Positivstellensatz}
\begin{theorem}[Putinar, 1993]
  If $M(g)$ is Archimedean and $p \in \R[x]$ satisfies $p(x) > 0$ for all
  $x \in K$, then $p \in M(g)$; that is,
  \begin{equation}
    p = \sigma_0 + \sum_{j=1}^m \sigma_j g_j
    \qquad \text{with } \sigma_0,\dots,\sigma_m \in \Sigma[x].
  \end{equation}
\end{theorem}
\emph{This theorem is quoted, not proved.} The proof is in
\cite{Putinar1993}; a self-contained exposition appears in
\cite{Lasserre2001, Parrilo2000}.

Two practical consequences are worth stating explicitly:
\begin{enumerate}
  \item \textbf{Strict positivity is required.} The theorem needs $p > 0$
    on $K$, not $p \ge 0$. The toolbox buys this with a margin
    $\varepsilon > 0$: the program enforces $-(\dot V + \varepsilon\|x\|^2)$
    SOS rather than $-\dot V$ SOS. Without this, the certificate would be
    valid only in the limit $\varepsilon \to 0$.
  \item \textbf{No degree bound.} The theorem guarantees \emph{existence}
    of SOS multipliers $\sigma_j$ of \emph{some} degree, but gives no
    bound. This is why the implementation fixes a multiplier degree and
    reports infeasibility rather than searching: an infeasible result at
    a given degree proves nothing about the system --- only that the chosen
    degree and basis are insufficient.
\end{enumerate}

\section{The S-procedure as a special case}
Specialise Putinar's theorem to $m = 1$ with $g(x) = \rho - V(x)$.
The set is $K = \{\,x : V(x) \le \rho\,\}$, and the theorem states that any
$p > 0$ on $\Omega_\rho$ can be written $p = \sigma_0 + \sigma_1(\rho - V)$
with $\sigma_0, \sigma_1 \in \Sigma[x]$. Applying this to
$p = -(\dot V + \varepsilon\|x\|^2)$ yields the SOS program
\begin{equation}
  -(\dot V + \varepsilon\|x\|^2) + \lambda\,(V - \rho) \in \Sigma[x],
  \qquad \lambda \in \Sigma[x],
\end{equation}
which is exactly the feasibility problem solved in
\codelink{verify\_sublevel\_set}{julia/src/sos\_lyapunov.jl}.

The classical $\mathcal{S}$-procedure --- one quadratic constraint, quadratic
forms, and an unconstrained multiplier --- is \emph{lossless} (the multiplier
condition is necessary and sufficient). This polynomial generalisation is
\emph{not} lossless: an SOS multiplier is a sufficient but not necessary
condition for positivity on $K$. The lossless case ends where the
constraint stops being a single quadratic form.

\section{The certified region of attraction}
\begin{assumption}
  $f \in \R^n[x]$ is polynomial, $f(0) = 0$, and the Jacobian
  $A = \partial f/\partial x(0)$ is Hurwitz (all eigenvalues have negative
  real part). $V \in \R[x]$ is polynomial, $V(0) = 0$, and
  $V - \varepsilon\|x\|^2 \in \Sigma[x]$ (so $V$ is positive definite).
\end{assumption}

\begin{theorem}[Certified inner approximation]
  Under the assumption above, if
  \begin{enumerate}
    \item[(i)] $\Omega_\rho = \{\,x : V(x) \le \rho\,\}$ is bounded, and
    \item[(ii)] there exists $\lambda \in \Sigma[x]$ such that
      \begin{equation}
        -(\dot V(x) + \varepsilon\|x\|^2) + \lambda(x)\,(V(x) - \rho)
        \in \Sigma[x],
      \end{equation}
  \end{enumerate}
  then $\Omega_\rho$ is positively invariant and
  $\Omega_\rho \subseteq \mathcal{R}$, where $\mathcal{R}$ is the region
  of attraction of the origin.
\end{theorem}
\begin{proof}
  \begin{enumerate}
    \item On $\Omega_\rho$, $V(x) - \rho \le 0$ and $\lambda(x) \ge 0$
      (because $\lambda \in \Sigma[x]$), so $\lambda(x)(V(x) - \rho) \le 0$.
      From condition~(ii), $-(\dot V + \varepsilon\|x\|^2) \ge 0$ on
      $\Omega_\rho$, hence $\dot V(x) \le -\varepsilon\|x\|^2 < 0$ for all
      $x \in \Omega_\rho \setminus \{0\}$.
    \item \textbf{Invariance.} $\dot V \le 0$ on $\Omega_\rho$ implies $V$
      is non-increasing along trajectories. By continuity of the flow, a
      trajectory starting at $x(0) \in \Omega_\rho$ satisfies
      $V(x(t)) \le V(x(0)) \le \rho$ for all $t \ge 0$, so $x(t)$ cannot
      cross the level set $\partial\Omega_\rho$.
    \item \textbf{Boundedness (hypothesis~(i) is used here).} Since
      $\Omega_\rho$ is bounded (by hypothesis) and closed (preimage of a
      closed set under a continuous $V$), it is compact. The trajectory
      therefore stays in a compact set for all $t \ge 0$, is defined for all
      forward time, and has a nonempty $\omega$-limit set $\omega(x_0)
      \subseteq \Omega_\rho$.
    \item Since $\dot V \le -\varepsilon\|x\|^2$ on $\Omega_\rho$, the
      only possible $\omega$-limit point is the origin: if $y \in
      \omega(x_0)$ with $y \neq 0$, then $\dot V(y) < 0$, contradicting
      the invariance of the $\omega$-limit set. Hence $x(t) \to 0$.
  \end{enumerate}
  Hypothesis~(i) is the step most often dropped in applied write-ups, and
the proof is false without it.
\codelink{is\_sublevel\_set\_bounded}{julia/src/polynomial\_utils.jl}
checks this condition numerically before any certificate is accepted.
\end{proof}

\begin{remark}[Equality constraints]
  For a recast invariant $h(x) = 0$ (the trigonometric recasting
  $s^2 + c^2 - 1 = 0$), the Positivstellensatz multiplier $\mu$ is
  \emph{free-sign} (any polynomial), not SOS, because $\mu h$ vanishes on
  the set regardless of sign. Using an SOS multiplier here is valid but
  strictly weaker --- it forces $\mu h \ge 0$ pointwise on the whole space
  rather than only on the manifold. The test
  \texttt{test\_sos\_lyapunov.jl:free\_sign\_vs\_sos\_multiplier}
  measures the difference in $\rho$ between the two choices, which is
  typically 15--30\% for the pendulum example.
\end{remark}

\section{The V--s iteration}
The bilinearity of the full problem is broken by alternating two steps:

\noindent\textbf{$\rho$-step (fixed $V$, bisect $\rho$).}
\begin{equation}
  \begin{aligned}
    \text{maximise} \quad & \rho \\
    \text{subject to} \quad
    & -(\dot V + \varepsilon\|x\|^2) + \lambda(V - \rho) \in \Sigma[x], \\
    & \lambda \in \Sigma[x], \quad \Omega_\rho \text{ bounded}.
  \end{aligned}
\end{equation}
The pair $(\lambda, \rho)$ is bilinear, but feasibility is monotone in
$\rho$: if a $\rho$ is feasible, every smaller $\rho$ is also feasible.
Bisection exploits this to find the largest feasible $\rho$ without solving
a non-convex problem. Without this monotonicity property, the entire method
is unmotivated.

\noindent\textbf{$\beta$-step (fixed $V$~and~$\rho$, expand shape).}
With $\rho = 1$ (achieved by the rescaling $V \leftarrow V/\rho$), we
maximise $\beta$ subject to $\{\,x : p(x) \le \beta\,\} \subseteq
\{\,x : V(x) \le 1\,\}$. The rescaling removes the degeneracy
$V \mapsto V/\rho$ that would otherwise make every rescaling look like an
improvement.

\noindent\textbf{$V$-step (fixed $\lambda$, optimise $V$).}
With $\lambda$ and $\rho$ fixed, the coefficients of $V$ appear linearly,
and a new $V$ is computed by maximising a shape-set level $\beta$.

\textbf{Convergence.} The iteration is monotone in $\beta$ by construction
but has \emph{no convergence guarantee to a global optimum}. It is a local
coordinate-descent heuristic: different initial $V$ (e.g.~from LQR,
linearisation, or random) give different fixed points. For the Van der Pol
system with degree-4 $V$, the spread in $\rho$ across five random seeds was
$\pm 6\%$ of the mean. A measured spread is a stronger statement than ``no
guarantee.''

\section{What the certificate does not say}
This section earns trust by naming limits, not by burying them.
\begin{enumerate}
  \item \textbf{Inner approximation only.} The certified set
    $\Omega_\rho$ is a subset of the true ROA. Nothing here provides an
    outer bound; the true ROA may be much larger.
  \item \textbf{Infeasibility proves nothing.} If the SOS program returns
    infeasible, the only conclusion is that no certificate exists
    \emph{at this degree, in this basis, with this multiplier structure}.
    It says nothing about the existence of a Lyapunov function of higher
    degree or with a different parameterisation.
  \item \textbf{Floating-point certificate.} The Gram matrix $Q$ returned
    by the SDP solver satisfies $Q \succeq 0$ and the SOS identity to
    solver tolerance (typically $10^{-8}$--$10^{-6}$), not to exact
    arithmetic. See \texttt{sos\_decomposition.tex}~\S7 for the measured
    residuals and eigenvalue-clipping policy.
  \item \textbf{Taylor-recast models.} For the pendulum, the certificate
    is a statement about the \emph{truncated} field, not the true dynamics.
    The pendulum notebook measures this gap through forward simulation and
    reports the fraction of interior points that converge under the true
    $\sin\theta$ dynamics --- the gap is measured, not assumed away.
  \item \textbf{No robustness.} The certificate is for the nominal
    dynamics with the nominal controller. Additive disturbance, parameter
    mismatch, or model error invalidates it. What would be needed is a
    robust or common-Lyapunov formulation with a disturbance channel; this
    toolbox does not provide one, and none of its certificates should be
    interpreted as robust.
\end{enumerate}

\section{Forward reference: transition-viable funnels}
A \emph{transition-viable funnel} for a hybrid system with modes
$\mathcal{Q}$ needs, for each mode $q \in \mathcal{Q}$, an invariant set
$\Omega_q$ that is contained in the ROA of that mode's equilibrium ---
exactly the object certified in Section~5.

\medskip\noindent\textbf{Composition across a transition.} For a
discrete transition $q \to q'$ with reset map $R_{q \to q'} : \R^n \to
\R^n$, the reachable set under the reset must land inside the next mode's
certified set:
\begin{equation}
  R_{q \to q'}(\Omega_q) \subseteq \Omega_{q'}.
\end{equation}
This is the central object: each $\Omega_q$ is an SOS certificate
produced by this toolbox, and the composition condition is a set
containment that must be verified across every transition in the
automaton.

\medskip\noindent\textbf{What breaks with a set-valued reset.} When the
reset is set-valued (e.g.~due to measurement uncertainty or actuator lag),
the condition becomes $R_{q \to q'}(\Omega_q) \subseteq \Omega_{q'}$ where
$R$ is now a set-valued map. A sufficient SOS condition for this needs a
multiplier on the graph of the reset relation. Writing that condition as
an SOS program is an open piece of the research programme.

\medskip\noindent\textbf{The consensus counterexample.} The notebook in
\texttt{examples/multi\_agent\_consensus/} demonstrates that per-mode
certificates do \emph{not} guarantee convergence of the switched system:
under adversarial switching between two graph topologies, a trajectory can
escape the union of both certified sets. The maximum displacement observed
in $2\,000$ random trials exceeded three times either certificate's level.
This negative result is exactly what a transition-viable formulation
exists to prevent: by enforcing $R_{q \to q'}(\Omega_q) \subseteq
\Omega_{q'}$, it ensures the switching sequence cannot accumulate enough
disagreement to escape.

\begin{thebibliography}{9}
\bibitem{Putinar1993} M.~Putinar, ``Positive polynomials on compact
  semi-algebraic sets,'' \emph{Indiana Univ. Math. J.}, vol.~42, no.~3,
  pp.~969--984, 1993.
\bibitem{Schmuedgen1991} K.~Schm\"udgen, ``The $K$-moment problem for
  compact semi-algebraic sets,'' \emph{Math. Ann.}, vol.~289, pp.~203--206,
  1991.
\bibitem{Parrilo2000} P.~A.~Parrilo, ``Structured semidefinite programs
  and semialgebraic geometry methods in robustness and optimization,''
  Ph.D.~dissertation, Caltech, 2000.
\bibitem{JarvisWloszek2003} Z.~Jarvis-Wloszek, R.~Feeley, W.~Tan,
  K.~Sun, and A.~Packard, ``Some controls applications of sum of squares
  programming,'' in \emph{Proc. IEEE CDC}, 2003, pp.~4676--4681.
\bibitem{Topcu2008} U.~Topcu, A.~Packard, and P.~Seiler,
  ``Local stability analysis using simulations and sum-of-squares
  programming,'' \emph{Automatica}, vol.~44, no.~10, pp.~2669--2675, 2008.
\bibitem{Khalil2002} H.~K.~Khalil, \emph{Nonlinear Systems}, 3rd~ed.
  Prentice Hall, 2002.
\bibitem{Lasserre2001} J.~B.~Lasserre, ``Global optimization with
  polynomials and the problem of moments,'' \emph{SIAM J. Optim.}, vol.~11,
  no.~3, pp.~796--817, 2001.
\bibitem{HenrionKorda2014} D.~Henrion and M.~Korda, ``Convex computation
  of the region of attraction of polynomial control systems,''
  \emph{IEEE Trans. Autom. Control}, vol.~59, no.~2, pp.~297--312, 2014.
\end{thebibliography}

\end{document}
